% !TeX program = xelatex
% UTF-8 transcription of the 19 photographs in chronological filename order.
% Source wording, abbreviations, numerical values and clinical claims are retained.
% Bold text represents yellow highlighting. Uncertain readings are marked explicitly.
% Standalone XeLaTeX/XePersian document; also usable with \input in a XePersian book.
\newif\ifNEStandalone
\ifdefined\settextfont
  \NEStandalonefalse
\else
  \NEStandalonetrue
\fi
\ifNEStandalone
  \documentclass[a4paper,12pt]{extbook}
  \usepackage[margin=2.2cm]{geometry}
  \usepackage{amsmath,amssymb,graphicx,array}
  \usepackage{xepersian}
  \settextfont{Amiri}
  \setlatintextfont{FreeSerif}
  \pagestyle{plain}
\fi
\ifNEStandalone\begin{document}\fi
\begingroup
\emergencystretch=3em
\let\NEolditemize\itemize
\renewcommand{\itemize}{\NEolditemize\setlength{\itemsep}{2pt}\setlength{\parsep}{0pt}\setlength{\topsep}{4pt}}
\let\NEoldenumerate\enumerate
\renewcommand{\enumerate}{\NEoldenumerate\setlength{\itemsep}{2pt}\setlength{\parsep}{0pt}\setlength{\topsep}{4pt}}
\newcommand{\NEen}[1]{\lr{#1}}
\newcommand{\NEmath}[1]{\lr{$#1$}}
\newcommand{\NEhl}[1]{\textbf{#1}}
\newcommand{\NEunclear}[1]{#1 [خوانش نامطمئن]}
\newif\ifNEFirstPage
\NEFirstPagetrue
\newcommand{\NEpage}[1]{\ifNEFirstPage\NEFirstPagefalse\else\clearpage\fi\section*{صفحهٔ #1}
  {\footnotesize\noindent\NEen{Date: / /}\quad
  \NEen{Sat. Sun. Mon. Tue. Thu. Wed. Fri.}\par
  \noindent\NEen{Subject:} ۱ ـ شرح حال و معاینهٔ نورولوژیک\par}\medskip}
\newcommand{\NEarrow}{\ensuremath{\longrightarrow}}
\newcommand{\NEphoto}[2]{%
  \IfFileExists{images/#1.jpg}{\includegraphics[#2]{images/#1.jpg}}%
  {\includegraphics[#2]{neurology/neurological examination/images/#1.jpg}}}

\chapter{شرح حال و معاینهٔ نورولوژیک}

% pic01.jpg = 20261012_215841925.jpg
\NEpage{۱}
\begin{latin}
\begin{itemize}
\item Confusion \NEarrow{} Memory loss, Difficulty understanding or producing spoken or written language, problems with numbers, faulty judgment, personality change.
\item \NEhl{Dizziness}:
  \begin{itemize}
  \item \NEhl{Vertigo} (illusion of movement of oneself or the environment).
  \item \NEhl{Imbalance} (Unsteadiness due to \NEhl{extrapyramidal}, \NEhl{vestibular}, \NEhl{cerebellar} or \NEhl{sensory} deficits).
  \item Presyncope (light-headedness resulting from cerebral hypoperfusion).
  \end{itemize}
\item Weakness \NEarrow{} Loss of power from disorders affecting motor pathways in the \NEhl{CNS} or \NEhl{PNS} or \NEhl{muscles}.
\item Shaking \NEarrow{} Tremor, chorea, athetosis, myoclonus, fasciculation.
\item Numbness \NEarrow{} hypesthesia (decreased sensitivity), hyperesthesia, \NEhl{paresthesia} (pins and needles sensation).
\item Spells \NEarrow{} Episodic and often recurrent symptoms \NEarrow{} Epilepsy, Syncope.
\item Quality of Symptoms:
  \begin{itemize}
  \item Neuropathic pain \NEarrow{} Direct injury to nerve.
  \item Described as unpleasant (dysesthetic), can be accompanied by hypersensitivity to pain (hyperalgesia) or touch (hyperesthesia) or perception of a normally stimulus as painful (\NEhl{allodynia}).
  \end{itemize}
\end{itemize}
\end{latin}
\NEhl{بارداری} یک فاکتور مستعدکننده برای \NEhl{سودوتومور سربری} است.

% pic02.jpg = 20261012_215848788.jpg
\NEpage{۲}
\begin{itemize}
\item \NEhl{سندرم ورنیکه ـ کورساکوف}: کمبود تیامین (\NEen{$B_1$}).
\item هایپرویتامینوز \NEen{A}: سودوتومور سربری (\NEhl{سردرد، ضعف بینایی، ادم پاپی}).
\item هایپرویتامینوز \NEen{$B_6$} (پیریدوکسین): \NEhl{پلی‌نوروپاتی}.
\end{itemize}
\subsection*{معاینهٔ سر و گردن}
\begin{itemize}
\item شکستگی قاعدهٔ جمجمه:
  \begin{itemize}
  \item هماتوم پشت اوریکولار (\NEen{battle sign}).
  \item هماتوم پری‌اوربیتال (\NEen{raccoon eyes}).
  \item هموتیمپان.
  \item اتوره، رینورهٔ \NEen{CSF}.
  \end{itemize}
\item مننژیت: برودزینسکی: فلکشن اندام تحتانی در مقابل فلکشن گردن.
\item \NEen{Kernig}: مقاومت در برابر اکستنشن زانو زمانی که هیپ فلکس است.
\item تست \NEen{SLR} (لازگ ساین): \NEen{$L_2$--$L_4$}، سیاتیک.
\item \NEunclear{\NEen{r-SLR}}: \NEhl{\NEen{$L_2$--$L_4$}}، فمورال \NEen{nerve}.
\end{itemize}
\subsection*{معاینهٔ نورولوژیک}
\begin{latin}
\NEhl{Mental status examination}
\begin{itemize}
\item Is level of consciousness normal or abnormal?
\item Is the cognitive function N/L?
\item If No \NEarrow{} Nature and extent.
\end{itemize}
\end{latin}

% pic03.jpg = 20261012_215853645.jpg
\NEpage{۳}
\subsection*{\NEen{A. N/L level of consciousness}}
\NEen{Awake, Alert, Oriented}.
\begin{latin}
\begin{center}
\footnotesize\setlength{\unitlength}{.75pt}
\begin{picture}(340,155)
\put(130,35){\vector(0,1){105}}
\put(130,35){\vector(1,0){190}}
\put(126,125){\line(1,0){4}}
\put(126,95){\line(1,0){4}}
\put(126,70){\line(1,0){4}}
\put(126,45){\line(1,0){4}}
\put(0,122){\makebox(120,12)[r]{Coherent}}
\put(0,92){\makebox(120,12)[r]{Purposeful}}
\put(0,67){\makebox(120,12)[r]{Semi-purposeful}}
\put(0,42){\makebox(120,12)[r]{Reflexive or none}}
\put(165,32){\line(0,1){6}}
\put(225,32){\line(0,1){6}}
\put(285,32){\line(0,1){6}}
\put(140,16){\makebox(50,12){Verbal}}
\put(140,3){\makebox(50,12){or visual}}
\put(203,16){\makebox(45,12){Tactile}}
\put(260,16){\makebox(50,12){Painful}}
\put(145,125){\circle{6}}
\put(151,122){\makebox(40,12)[l]{N/L}}
\put(285,48){\circle{6}}
\put(291,48){\makebox(45,12)[l]{Coma}}
\end{picture}
\end{center}
\end{latin}
\subsection*{\NEen{B. Cognitive function}}
هدف این است که بررسی کنیم آیا اختلالی در عملکرد شناختی فرد وجود دارد یا نه و اگر وجود دارد مربوط به ناحیهٔ خاصی از مغز است یا خیر.
\begin{itemize}
\item \NEhl{\NEen{Aphasia}} (\NEen{disorders of language})، \NEhl{\NEen{Memory}} (\NEen{amnesia}): مکان خاص دارند.
\end{itemize}
\subsection*{\NEhl{\NEen{1. Bifrontal or diffuse functions}}}
\begin{itemize}
\item \NEen{Attention, Concentration}:
  \begin{itemize}
  \item \NEen{A}: بلا‌فاصله پس از ارائهٔ اعداد، آنها را تکرار کند؛ شامل ۷ تا.
  \item \NEen{C}: از ۱۰۰، ۷ تا ۷ تا کم کند.
  \end{itemize}
\item \NEen{Abstract thought processes}: \NEen{Insight, Judgment}.
\item \NEen{Fund of knowledge}: ستاره‌های ورزشی، رئیس‌جمهور فعلی.
\item \NEen{Affect}: \NEen{talkativeness, facial expression, posture}؛ غم و شادی خود.
\item \NEen{Delusions, hallucinations}: \NEen{Psychiatric disease, Confusional states}؛ \NEen{Alcohol withdrawal}.
\end{itemize}
\subsection*{\NEhl{\NEen{2. Memory}}}
\begin{itemize}
\item \NEen{Diffuse cortical or bilateral temporal}.
\item \NEen{Immediate recall, Recent memory, Remote memory}.
\item \NEen{Immediate recall}: ارزیابی همان \NEen{attention}.
\item \NEen{Recent memory}: ۳–۵ دقیقه بعد.
\item \NEen{Confusional states}: معمولاً \NEen{immediate} مختل می‌شود.
\item \NEen{Amnesia}: معمولاً \NEen{recent} مختل است.
\item ناتوانی در به خاطر آوردن اسم خود در فرد هوشیار: اختلال روان‌پزشکی است.
\end{itemize}

% pic04.jpg = 20261012_215900463.jpg
\NEpage{۴}
\subsection*{\NEhl{\NEen{3. Language}}}
\NEen{Comprehension, Repetition, Fluency, Naming, Reading, Writing}.
\begin{latin}
\begin{center}
\begin{tabular}{lccc}
 & Fluency & Comprehension & Repetition\\\hline
Expressive (Broca) & $-$ & $+$ & $-$\\
Receptive (Wernicke) & $+$ & $-$ & $-$\\
Global & $-$ & $-$ & $-$\\
Conduction & $+$ & $+$ & $-$\\
Transcortical Expressive & $-$ & $+$ & $+$\\
Transcortical Receptive & $+$ & $-$ & $+$\\
Transcortical Global & $-$ & $-$ & $+$\\
Anomic (naming) & $+$ & $+$ & $+$\\
\end{tabular}
\end{center}
\end{latin}
% The original numbered connectivity drawing is retained without a new image file.
\begin{center}
\NEphoto{pic04}{width=.85\linewidth,trim=300bp 862bp 427bp 1610bp,clip}
\end{center}
\begin{latin}
\begin{enumerate}
\item Motor Speech area: Broca.
\item Language Comprehension area: Wernicke.
\item Arcuate Fasciculus: Conductive aphasia.
\item Transcortical Expressive aphasia.
\item Transcortical Receptive aphasia.
\end{enumerate}
\end{latin}

% pic05.jpg = 20261012_215907345.jpg
\NEpage{۵}
\subsection*{\NEen{Broca aphasia}}
\begin{itemize}
\item \NEen{Agrammatical, telegraphic nature of the little speech}.
\item \NEen{Unable to write or repeat normally}.
\item دستورات را اجرا می‌کند و متوجه می‌شود.
\item نسبت به شرایط آگاه است، احساس بدی دارد.
\end{itemize}
\subsection*{\NEen{Wernicke aphasia}}
\begin{itemize}
\item \NEen{Paraphasic error}: واژه‌های آوایی مشابه.
\item \NEen{Neologisms}.
\item نوشتن نیز مثل تکلم کردن مختل است.
\item نمی‌تواند فرمان‌های صدایی یا نوشتاری را پیروی کند.
\item اگر ادا در بیاوریم می‌تواند تقلید کند.
\item نسبت به شرایطشان ناآگاه هستند.
\end{itemize}
\subsection*{\NEen{Speech disorder (Dysarthria)}}
متفاوت از زبان، \NEunclear{معنی کلمه هستش} و \NEunclear{کلامش را تغییر نمی‌کند}.
% Preserve the two uncertain handwritten lines for direct comparison.
\begin{center}
\NEphoto{pic05}{width=.65\linewidth,trim=850bp 1162bp 82bp 1350bp,clip}
\end{center}
\subsection*{\NEhl{\NEen{4. Sensory Integration}}}
ناحیهٔ درگیری: \NEhl{لوب پریتال}.
\begin{itemize}
\item \NEen{Misperception of or inattention to sensory stimuli}.
\item در سمت مخالف ضایعه.
\item \NEhl{\NEen{Astereognosis}}: \NEen{Inability to identify objects by touch placed in the hand}.
\item \NEhl{\NEen{Agraphesthesia}}: ناتوانی در فهمیدن عددی که روی پوست نوشته شود.
\item \NEen{Failure of two point discrimination}.
\item \NEhl{\NEen{Allesthesia}}: \NEen{Misplaced localization of a tactile stimulus}؛ \NEen{Typically more proximal}.
\end{itemize}

% pic06.jpg = 20261012_215914185.jpg
\NEpage{۶}
\begin{itemize}
\item \NEen{Extinction}: \NEen{Failure to perceive a visual or tactile stimulus while applied bilaterally, even though it can be perceived when applied unilaterally}.
\item \NEen{Neglect}: \NEen{Failure to attend to space or use the limbs on one side of the body}.
  \begin{itemize}
  \item بیمار با ضایعه در \NEhl{پریتال \NEen{R}}: نگلکت سمت مقابل \NEen{L}.
  \item اگر بگوییم ساعت بکشد، همهٔ اعداد را در نیمهٔ \NEhl{\NEen{R}} قرار می‌دهد.
  \end{itemize}
\item \NEhl{\NEen{Anosognosia}}: \NEen{Unawareness of a neurologic deficit}.
\item \NEen{Constructional Apraxia}: \NEen{Inability to draw accurate representations of external space}.
\end{itemize}
\subsection*{\NEhl{\NEen{5. Motor Integration}}}
\begin{itemize}
\item \NEen{Apraxia}: \NEen{Inability to perform previously learned tasks despite intact motor and sensory function}.
\item \NEen{Ask patient to simulate the use of a key, comb or fork}.
\item یک نوع اختلال آپراکسی: \NEhl{\NEen{Contralateral premotor frontal}} \NEen{cortex lesion}.
\item سایر اختلال آپراکسی: مثل آپراکسی گیت: \NEen{bifrontal or diffuse lesions}.
\end{itemize}
\subsection*{\NEhl{\NEen{Cranial Nerves}}}
عصب اپتیک (\NEen{II}):
\begin{itemize}
\item افتالموسکوپی مستقیم.
\item حدت بینایی: \NEen{Measuring Visual Acuity}.
\item \NEen{Mapping the visual fields}.
\end{itemize}

% pic07.jpg = 20261012_215919442.jpg
\NEpage{۷}
یادداشت بالای صفحه: داروهای میدریاتیک = سمپاتومیمتیک + آنتی‌کولینرژیک.
\subsection*{۱ ـ افتالموسکوپی}
باید در یک اتاق تاریک باشد تا مردمک‌ها دیلاته شوند.
\begin{itemize}
\item اپتیک دیسک نرمال:
  \begin{itemize}
  \item \NEen{yellowish, oval}.
  \item \NEhl{مارژین شارپ}.
  \item \NEhl{پالس‌های خودبه‌خودی در وریدها} دیده می‌شود.
  \end{itemize}
\item \NEen{papilledema}: تورم دیسک اپتیک ناشی از \NEmath{ICP\uparrow}.
  \begin{itemize}
  \item وریدها پالس ندارند و متسع شده‌اند.
  \item خونریزی خطی در بوردر دیسک اپتیک.
  \item مارژین دیسک محو می‌شود؛ شروع از سمت نازال.
  \item تقریباً همواره دوطرفه است.
  \item معمولاً بینایی را مختل نمی‌کند؛ فقط نقطهٔ کور بزرگ‌تر می‌شود.
  \end{itemize}
\item \NEen{Optic disk pallor}: \NEen{Atrophy of the optic nerve}.
  \begin{itemize}
  \item \NEhl{\NEen{MS patients}}.
  \item \NEen{Associated with}: \NEen{defects in visual acuity, visual fields, pupillary reactivity}.
  \end{itemize}
\end{itemize}
\subsection*{\NEen{2. Visual acuity}}
\begin{itemize}
\item هر چشم جداگانه معاینه.
\item \NEen{Snellen Eye chart (6m), Rosenbaum pocket eye chart (36cm)}.
\item \NEen{N/L: 20/20}.
\item نقص‌های قرنیه ـ لنز ـ رتین ـ آسیب عصب اپتیک، به‌صورت غیرمساوی دچار اختلال می‌شوند.
\end{itemize}
\subsection*{\NEen{3. Visual fields}}
\begin{itemize}
\item معاینهٔ هر چشم جداگانه.
\item \NEen{Confrontation technique}: با فاصلهٔ یک بازوی دست از بیمار.
\item \NEen{Central Scotoma}: نقطهٔ کور؛ در نیمهٔ تمپورال میدان بینایی.
\end{itemize}

% pic08.jpg = 20261012_215924530.jpg
\NEpage{۸}
\begin{center}
\NEphoto{pic08}{width=.68\linewidth,trim=320bp 890bp 320bp 260bp,clip}
\end{center}
\begin{itemize}
\item برچسب‌های شکل مسیر بینایی: \NEen{L, R, T, N, optic nerve, optic Tract, optic radiation}؛ هستهٔ ژنیکولیت جانبی.
\item محل‌های شماره‌گذاری‌شده در شکل: ۱ در انتهای خلفی چشم چپ؛ ۲ روی عصب اپتیک چپ؛ ۳ در وسط کیاسما؛ ۴ روی بخش خارجی عصب اپتیک راست؛ ۵ روی اپتیک تراکت چپ؛ ۶ و ۷ روی اپتیک رادییشن چپ.
\item ۱ ـ \NEen{Central Scotoma}: در شکل، نقطهٔ مرکزی میدان \NEen{L} تیره و میدان \NEen{R} خالی است.
  \begin{itemize}
  \item التهاب دیسک اپتیک (\NEhl{نوریت اپتیک}).
  \item التهاب عصب اپتیک (نوریت رتروبولبار).
  \end{itemize}
\item ۲ ـ نابینایی کامل چشم چپ.
\end{itemize}

% pic09.jpg = 20261012_215930685.jpg
\NEpage{۹}
% Retain the seven paired visual-field drawings directly from pic09.
\begin{center}
\NEphoto{pic09}{height=.60\textheight,trim=120bp 380bp 1140bp 280bp,clip}
\end{center}
\begin{enumerate}
\setcounter{enumi}{2}
\item \NEen{Bitemporal hemianopia}: فشار به \NEhl{کیاسمای بینایی}؛ تودهٔ هیپوفیز.
\item \NEen{Right nasal hemianopia}: \NEhl{اینترنال کاروتید آنوریسم}.
\item \NEhl{\NEen{Right homonymous hemianopia}}: آسیب اپتیک تراکت چپ.
\item \NEen{Right homonymous superior quadrantanopia}: آسیب به لوب تمپورال چپ (\NEhl{\NEen{Meyer Loop}}).
\item \NEen{Right homonymous inferior quadrantanopia}: آسیب به \NEhl{پریتال چپ}.
\item \NEen{Right homonymous hemianopia}: آسیب کامل \NEen{Left optic radiation}.
\item \NEen{Right homonymous hemianopia with macular sparing}: آسیب \NEhl{پوستریور سربرال آرتری اکلوژن}.
\end{enumerate}
\NEen{L} و \NEen{R}: برچسب‌های هر جفت میدان بینایی در شکل.

% pic10.jpg = 20261012_215935611.jpg
\NEpage{۱۰}
\subsection*{\NEen{C. Oculomotor (III), Trochlear (IV), Abducens (VI) Nerves}}
\begin{latin}
\begin{center}\small
\begin{tabular}{p{.16\linewidth}p{.23\linewidth}p{.15\linewidth}p{.10\linewidth}p{.19\linewidth}}
Name & Appearance & Light reaction & Near reaction & Lesion site\\\hline
Adie (tonic) pupil & Unilateral large pupil & Sluggish & N/L & Ciliary ganglion\\
Argyll Robertson pupil & Bilateral small, irregular pupils & \NEhl{Absent} & N/L & \NEhl{Midbrain}\\
Horner Syndrome & Unilateral small pupil and ptosis & Normal & N/L & Sympathetic innervation of eye\\
Marcus Gunn pupil & Normal & Consensual $>$ direct & N/L & \NEhl{Optic nerve}\\
\end{tabular}
\end{center}
\end{latin}
نام؛ ظاهر؛ پاسخ مردمک به نور؛ پاسخ مردمک؛ محل ضایعه.
\subsection*{\NEen{1. pupils}}
معاینه: شکل، اندازه‌اش، نور، نزدیکی.
\begin{itemize}
\item \NEen{N/L}: \NEhl{\NEmath{\sim3\ \mathrm{mm}}} / \NEen{6mm in children} / \NEen{2mm in the elderly}.
\item دو مردمک می‌توانند تا \NEen{1mm} اختلاف سایز داشته باشند: فیزیولوژیک آنیزوکوریا.
\item طی تطابق (\NEen{accommodation}) مردمک‌ها تنگ می‌شوند.
\item میوز: از طریق \NEhl{پاراسمپاتیک} که از \NEhl{\NEen{midbrain}} منشأ می‌گیرد و با عصب \NEhl{اکولوموتور} همراه می‌شود.
\item آسیب این مسیر: مردمک دیلاته (\NEunclear{\NEmath{>7\ \mathrm{mm}}})؛ \NEhl{غیر \NEen{reactive}}.
\item میدریاز: از طریق یک \NEhl{راه سمپاتیک ۳ نورونه}:
  \begin{itemize}
  \item هیپوتالاموس / \NEen{brainstem} / \NEen{$T_1$} / سوپریور سرویکال گانگلیون / چشم.
  \end{itemize}
\item آسیب این مسیر: مردمک \NEhl{میوتیک} (\NEunclear{\NEmath{<1\ \mathrm{mm}}})؛ \NEhl{\NEen{non-reactive}}.
\end{itemize}
\subsection*{\NEen{2. Eyelids and orbits}}
\begin{itemize}
\item اینترپالپبرال فیشر معمولاً حدود \NEen{10mm} است.
\item پلک فوقانی معمولاً حدود \NEen{1--2mm} از \NEen{Iris} را می‌پوشاند.
\end{itemize}

% pic11.jpg = 20261012_215940730.jpg
\NEpage{۱۱}
\begin{itemize}
\item پتوز: آسیب به عضلهٔ \NEhl{لواتور پالپبرا}.
  \begin{itemize}
  \item آسیب به عصب \NEhl{اکولوموتور (\NEen{III}) یا سمپاتیک مربوط به عضله}.
  \end{itemize}
\item سندرم هورنر: پتوز + میوز + انهیدروز همان.
\item بیرون‌زدگی چشم: اگزوفتالمی، پروپتوز.
\item معاینه از بالای سر و پشت بیمار.
\end{itemize}
\subsection*{\NEen{3. Eye movements}}
\NEhl{افتالموپلژی}: وقتی عملکرد یک عضلهٔ چشم مختل است، حرکات چشم در جهت آن عضله مختل است؛ انحراف به سمت مقابل پیدا می‌کند.
\begin{center}
\NEphoto{pic11}{width=.67\linewidth,trim=680bp 1110bp 320bp 1290bp,clip}
\end{center}
برچسب‌های شکل:
\begin{itemize}
\item چپ: لترال رکتوس؛ راست: مدیال رکتوس.
\item بالا در ابداکشن: سوپریور رکتوس؛ پایین در ابداکشن: اینفریور رکتوس.
\item بالا در اداکشن: اینفریور ابلیک؛ پایین در اداکشن: سوپریور ابلیک.
\end{itemize}
\NEen{Adducted eye is elevated by inferior oblique}.
\begin{itemize}
\item عصب‌دهی: سوپریور ابلیک: \NEhl{تروکلئار (\NEen{IV})}.
\item لترال رکتوس: \NEhl{ابدوسنس (\NEen{VI})}.
\item بقیه: اکولوموتور (\NEen{III}).
\item اگر بیمار دیپلوپی دارد باید بخواهیم به جهتی که بیشترین دوبینی را دارد چشم‌ها را حرکت دهد تا عضلهٔ درگیر را ببینیم. سپس چشم‌ها را به نوبت ببندد و بگوید کدام تصویر (دور یا نزدیک) پنهان می‌ماند.
\end{itemize}
\begin{latin}
The picture displaced farther in the direction of gaze is always referable to the weak eye.
\end{latin}

% pic12.jpg = 20261012_215946503.jpg
\NEpage{۱۲}
در ضعف لترال رکتوس چپ: دیپلوپی در \NEen{Leftward gaze} حداکثر است. وقتی چشم چپ بسته شود، تصویر خارجی‌تر محو می‌شود.
\subsection*{\NEhl{\NEen{4. Ocular oscillations}}}
\subsubsection*{\NEen{Nystagmus}}
\begin{itemize}
\item \NEhl{در انتهای \NEen{gaze} ارادی طبیعی است}.
\item داروهای ضدتشنج و آرام‌بخش.
\item درگیری عضلات خارج چشمی.
\item درگیری مسیرهای \NEhl{وستیبولار} یا \NEhl{سربلار / مخچه‌ای}.
\item شایع‌ترین نوع: \NEhl{\NEen{jerk nystagmus}}.
  \begin{itemize}
  \item \NEen{Slow phase of movement followed by a fast phase in the opposite direction}.
  \end{itemize}
\item معاینه: چشم‌ها در حالت عادی و سپس در هر ۴ \NEen{gaze} کاردینال مشاهده.
\item جهت، در کدام \NEen{gaze}، \NEen{amplitude} (فاین، کورس)، عوامل تشدیدکننده مثل \NEunclear{سرگیجه}، جهت سر.
\item نوع کمتر شایع: \NEen{pendular nystagmus}.
  \begin{itemize}
  \item معمولاً در شیرخواران شروع می‌شود / \NEhl{دو فاز آن معمولاً مساوی هستند}.
  \end{itemize}
\item نیستاگموس در پوزیشن اولیه با اختلال بینایی پاتولوژیک است: \NEen{end point}.
\item جزء سریع: \NEunclear{وسط‌تر} در پاتولوژیک است.
\end{itemize}
\subsection*{\NEen{D. Trigeminal Nerve (V)}}
شاخه‌ها:
\begin{itemize}
\item افتالمیک (\NEen{$V_1$}) (پیشانی).
\item ماگزیلاری (\NEen{$V_2$}) (گونه).
\item مندیبولار (\NEen{$V_3$}) (فک).
\end{itemize}

% pic13.jpg = 20261012_215952305.jpg
\NEpage{۱۳}
\subsection*{\NEen{Corneal Reflex}}
\begin{itemize}
\item پنبه روی قرنیه کشیده شود ـ اسکلرا سفید.
\item قوس رفلکس:
  \begin{itemize}
  \item \NEhl{حسی: تری‌ژمینال (\NEen{$V_1$})}.
  \item \NEhl{حرکتی: فیشیال (\NEen{VII})}.
  \end{itemize}
\item \NEen{N/L}: پلک زدن دوطرفه.
\item اگر هیچ‌کدام پلک نزدند: آسیب \NEen{$V_1$}.
\item فقط یک چشم پلک بزند: آسیب \NEen{VII} همان سمت نزده.
\end{itemize}
\subsection*{بررسی \NEen{motor} عصب تری‌ژمینال}
\begin{itemize}
\item \NEhl{قدرتش باز کردن و بستن دهان}.
\item فشار دادن دندان‌ها به هم. [در متن یک عبارت پیش از این خط خورده است: \NEunclear{فک را باز کردن}.]
\item بیمار دندان‌ها را فشار می‌دهد؛ ما نباید بتوانیم دهانش را باز کنیم.
\end{itemize}
\subsection*{\NEen{E. Facial (VII) Nerve}}
\begin{itemize}
\item عضلات صورت و \NEhl{حس چشایی \NEmath{\tfrac{2}{3}} قدامی زبان}.
\item \NEhl{قرینگی فیشر پالپبرال} دوطرف / قرینگی چین نازولبیال.
\item موتور فوقانی از \NEhl{هر دو نیمکره} / نیمهٔ صورت (چشم پایین) از \NEhl{نیمکرهٔ مقابل}.
\item حس زبان:
  \begin{itemize}
  \item لمسی: قدام: \NEen{V} / خلف: \NEhl{\NEen{IX}}.
  \item چشایی: قدام: \NEen{VII} / خلف: \NEhl{\NEen{IX}}.
  \end{itemize}
\item پیشانی را بالا ببرد / چشم‌ها را فشار دهد / \NEhl{خندیدن} / نشان دادن دندان‌ها.
\item \NEhl{درگیری سنترال \NEen{VII}: پیشانی مشکل ندارد}.
\item فشار دادن لب‌ها به یکدیگر / \NEhl{باد کردن لب‌ها}.
\item \NEhl{\NEen{Facial weakness}} معمولاً با \NEen{dysarthria} همراه است؛ در \NEhl{تلفظ \NEen{M}}.
\end{itemize}

% pic14.jpg = 20261012_215958132.jpg
\NEpage{۱۴}
\subsection*{\NEen{F. Acoustic (VIII) Nerve}}
\begin{itemize}
\item انجام \NEhl{اتوسکوپی} برای بررسی کانال و \NEen{TM}.
\item انجام تست \NEhl{وبر و رینه} با یک دیاپازون \NEen{512Hz}.
\item بررسی \NEhl{حدت شنوایی} با مالیدن دو انگشت خود در فاصلهٔ \NEen{5cm} از گوش بیمار.
\item اگر نشنید باید نوع ناشنوایی را مشخص شود.
\end{itemize}
\subsubsection*{تست رینه}
\begin{itemize}
\item بیس دیاپازون مرتعش روی \NEhl{ماستوئید} بیمار، و از آن زمانی که بیمار دیگر نمی‌شنود می‌پرسیم.
\item سپس نوک دیاپازون مرتعش را \NEhl{نزدیک مجرای خارجی گوش} می‌گیریم.
\item \NEen{N/L}: هدایت \NEmath{air > bone}.
\item سنسوری‌نورال: \NEmath{air > bone}.
\item کنداکتیو: \NEmath{bone > air}.
\end{itemize}
\subsubsection*{تست وبر}
\begin{itemize}
\item بیس دیاپازون مرتعش روی وسط پیشانی گذاشته شود.
\item \NEen{N/L}: میدلاین.
\item سنسوری‌نورال: در گوش سالم بهتر می‌شنود.
\item کنداکتیو: در گوش درگیر بهتر می‌شنود.
\end{itemize}
\subsubsection*{در بیمار با شکایت سرگیجهٔ وضعیتی}
\NEen{Nylen-Barany or} \NEhl{\NEen{Dix-Hallpike maneuver}}.
\begin{itemize}
\item برای ایجاد شرایط مشابه‌ای که منجر به سرگیجه.
\item بیمار روی تخت نشسته ـ روبه‌جلو؛ سپس سریع \NEen{Supine} و سر ۴۵ درجه خم بر لبهٔ تخت.
\item سپس دوباره نشسته، سر ۴۵ درجه به راست و ... .
\item سپس دوباره نشسته، سر ۴۵ درجه به چپ و ... .
\item چشم‌ها برای نیستاگموس و بیمار برای سرگیجه.
\end{itemize}

% pic15.jpg = 20261012_220002837.jpg
\NEpage{۱۵}
\subsection*{\NEen{G. Glossopharyngeal (IX) and Vagus (X) Nerves}}
\begin{itemize}
\item بررسی \NEhl{موتور} این دو عصب: از مریض می‌خواهیم بگوید \NEhl{\NEen{``ah''}} با دهان باز.
\item سپس نگاه می‌کنیم که \NEhl{آیا کام به‌صورت کامل و قرینه \NEen{elevate} شده است}.
\item در صورت ضعف عضلات کام: \NEen{dysarthria}: \NEhl{حرف \NEen{K}}.
\item بررسی \NEhl{سنسوری}: \NEhl{رفلکس \NEen{gag}}؛ انتهای زبان در دوطرف بررسی می‌شود. رفلکس باید مشابه و قرینه باشد.
\end{itemize}
\subsection*{\NEen{H. Accessory (XI) Nerve}}
\begin{itemize}
\item عضلهٔ \NEhl{\NEen{SCM}} و \NEhl{تراپزیوس} را عصب‌دهی می‌کند.
\item دست روی فک؛ بیمار سر خود را در مقابل مقاومت معاینه‌کننده بچرخاند: \NEen{SCM}.
\item بیمار شانه‌های خود را در مقابل مقاومت بالا ببرد و غیرقرینگی حس کنیم: تراپزیوس.
\end{itemize}
\subsection*{\NEen{I. Hypoglossal (XII) Nerve}}
\begin{itemize}
\item \NEhl{عضلات زبان} را عصب‌دهی می‌کند.
\item با \NEhl{زبان} لپ را به بیرون فشار دهد، در حالی که ما از خارج مقاومت کنیم.
\item \NEhl{خارج کردن زبان}: به سمت \NEhl{ضعیف} منحرف می‌شود.
\item ضعف عضلات زبان: \NEen{Dysarthria}: \NEhl{حرف \NEen{L}}.
\item اگر عصب قطع شود: آتروفی و فاسیکولیشن زبان.
\end{itemize}

% pic16.jpg = 20261012_220010196.jpg
\NEpage{۱۶}
\subsection*{\NEhl{\NEen{Motor Function}}}
\begin{latin}
\begin{itemize}
\item LMN: Weakness, $\downarrow$ muscle tone, $\downarrow$ reflex, atrophy, fasciculation.
\item UMN: Weakness, $\uparrow$ muscle tone, Upward \NEhl{Babinski}, $\uparrow$ reflex.
\end{itemize}
\end{latin}
\subsection*{\NEhl{\NEen{A. Bulk}}}
\begin{itemize}
\item آتروفی معمولاً ناشی از \NEhl{ضایعات \NEen{LMN}} است.
\item مقایسهٔ عضلات دوطرف برای ارزیابی آتروفی \NEhl{غیرقرینه}.
\item اندازه‌گیری قطر با متر نواری.
\item آتروفی می‌تواند همراه با فاسیکولیشن باشد.
\end{itemize}
\subsection*{\NEhl{\NEen{B. Tone}}}
\begin{itemize}
\item مقاومت عضله در برابر حرکات \NEhl{\NEen{passive}} در مفصل.
\item \NEen{N/L}: مقاومت اندکی وجود دارد.
\item هایپوتونی یا \NEen{flaccidity}: درگیری \NEhl{عضله، \NEen{LMN}، مخچه}.
\item هایپرتونی:
  \begin{itemize}
  \item \NEen{rigidity}: مقاومت در تمام \NEen{ROM} مفصل وجود دارد.
  \item \NEen{spasticity}: مقاومت در طول \NEen{ROM} متفاوت است.
  \end{itemize}
\item ریجیدیتی: \NEhl{بازال گانگلیا}.
\item اسپاستیسیتی: \NEhl{مسیر کورتیکواسپاینال}.
\item در مفصل آرنج معاینه می‌شود و مچ (در حالت طبیعی در ۹۰ باشد).
\item در پا: بیمار \NEen{Supine}، یک دست زیر زانو؛ سپس ناگهان بالا آورده شود.
\item در \NEen{N/L} پاشنهٔ پا نباید کامل از تخت جدا شود.
\item آنکل: سر پا به‌صورت \NEen{passive} چرخانده شود؛ نباید مقاومت کند در حالت \NEen{N/L}.
\item سر پا ـ جلو و عقب خم کنیم؛ نباید مقاومتی احساس کنیم.
\end{itemize}

% pic17.jpg = 20261012_220015400.jpg
\NEpage{۱۷}
\subsection*{\NEen{C. Strength}}
\begin{itemize}
\item مقیاس صفر تا پنج:
  \begin{itemize}
  \item ۵: نرمال.
  \item ۴: کاهش‌یافته ولی در مقابل جاذبه و مقاومت غلبه می‌کند.
  \item ۳: غلبه بر جاذبه ولی مقاومت نه.
  \item ۲: صرفاً حرکت افقی.
  \item ۱: فاسیکولیشن و حرکات جزئی.
  \item ۰: بدون حرکت.
  \end{itemize}
\item باید نیروی مقاومت ما تدریجی اضافه شود.
\item هدف اصلی: \NEhl{یافتن اختلاف} بین دوطرف راست و چپ اندام فوقانی و تحتانی.
\item \NEhl{\NEen{Pyramidal weakness}} (ضایعه در مسیر کورتیکواسپاینال):
  \begin{itemize}
  \item ضعف ترجیحی در \NEhl{اکستنسور و ابداکتور اندام فوقانی}.
  \item \NEhl{فلکسور در اندام تحتانی}.
  \item \NEen{Fine finger movements are slowed}: زدن انگشت اشاره به شست.
  \item \NEhl{\NEen{Pronator drift} وجود دارد}.
  \end{itemize}
\item با اختلال بیشتر \NEhl{دیستال ویکنس}: کاراکتریستیک \NEhl{پلی‌نوروپاتی} است.
\item با اختلال بیشتر \NEhl{پروگزیمال ویکنس}: \NEhl{میوپاتی}.
\end{itemize}
\subsection*{\NEhl{\NEen{Sensory Function}}}
\begin{itemize}
\item حس سوماتیک از طریق ستون خلفی نخاع به تالاموس می‌رود (لارج فایبر)؛ به مدیال لمنیسکوس برین‌استم.
\item و اسپاینوتالامیک تراکت (اسمال فایبر).
\item \NEhl{\NEen{pain and temperature}}: \NEhl{\NEen{Small Fibers}}. [واژهٔ \NEen{vibration} در این خط خط خورده و \NEen{Temperature} زیر آن جایگزین شده است.]
\item \NEen{light touch}: \NEen{both}.
\item \NEhl{\NEen{Vibration and position}}: \NEhl{\NEen{Large Fibers}}.
\item بررسی حسی از دیستال به پروگزیمال انجام می‌شود.
\end{itemize}

% pic18.jpg = 20261012_220021430.jpg
\NEpage{۱۸}
\subsection*{\NEen{C.} پوزیشن}
حرکت دادن انگشت پا، بیمار باید بگوید به کدام جهت.
\begin{itemize}
\item با چشم بسته نوک دو انگشت را به هم بزند.
\end{itemize}
\subsection*{\NEhl{\NEen{Coordination}}}
\begin{itemize}
\item \NEhl{آتاکسی یعنی کوردینیشن مختل}: معمولاً ضایعهٔ مخچه.
\item حرکات چشم، تکلم، اندام‌ها، تنه.
\item تست‌های عملکرد مخچه:
  \begin{itemize}
  \item \NEhl{\NEen{finger to nose}} (چشم باز).
  \item \NEhl{\NEen{Rebound}} (چشم بسته).
  \item \NEhl{\NEen{heel-knee-shin}} (چشم باز).
  \end{itemize}
\end{itemize}
\subsubsection*{\NEen{A. Limb ataxia}}
\begin{itemize}
\item زدن کف یا پشت دست به دست مقابل: بررسی دیس‌دیادوکوکینزی.
\item \NEen{finger-to-nose}.
\item آتاکسی ممکن است همراه با \NEhl{\NEen{Intention tremor}} باشد.
\item در ابتدا و انتهای هر حرکت واضح‌تر است.
\end{itemize}
\subsubsection*{\NEen{B. Truncal ataxia}}
روی لبهٔ تخت یا روی صندلی باید بتواند صاف بنشیند.
\subsection*{\NEhl{\NEen{Reflexes}}}
\subsubsection*{\NEen{A. Tendon reflexes}}
\begin{itemize}
\item \NEhl{پلی‌نوروپاتی: کاهش یا حذف رفلکس}.
\item ضایعات \NEhl{کورتیکواسپاینال: افزایش رفلکس}.
\item درجه‌بندی:
  \begin{itemize}
  \item \NEen{4. Very brisk often with clonus}.
  \item \NEen{3. brisk but N/L}.
  \item \NEen{2. N/L}.
  \item \NEen{1. Minimal}.
  \item \NEen{0. Absent}.
  \end{itemize}
\item هدف اصلی بررسی \NEhl{غیرقرینگی یا نبود رفلکس} است.
\end{itemize}

% pic19.jpg = 20261012_220025195.jpg
\NEpage{۱۹}
\begin{itemize}
\item نبودن سیمتریک رفلکس: پلی‌نوروپاتی.
\item افزایش سیمتریک رفلکس: ضایعهٔ دوطرفهٔ مغز یا طناب نخاعی.
\item \NEen{biceps and brachioradialis ($C_5$--$C_6$) / triceps ($C_7$--$C_8$)}.
\item \NEen{quadriceps ($L_3$--$L_4$) / Achilles ($S_1$--$S_2$)}.
\end{itemize}
\subsection*{\NEen{B.} \NEhl{\NEen{Superficial reflexes}}}
\begin{itemize}
\item در ضایعات \NEhl{مسیر کورتیکواسپاینال} تغییر می‌یابد.
\item \NEen{Plantar reflex}: انگشت پا به بالا آمده: \NEen{Babinski sign}.
\item رفلکس‌های دوران نوزادی و شیرخوارگی طبیعی است؛ به دنبال افزایش سن و یا ضایعات \NEhl{لوب فرونتال} دوباره ظاهر می‌شوند.
\item \NEen{palmar grasp / plantar grasp / palmomental reflex}.
\end{itemize}
\subsection*{\NEhl{\NEen{Stance and Gait}}}
\begin{latin}
\begin{itemize}
\item Romberg sign.
\item Walking normally, on the heels, in tandem.
\item Classic gaits:
  \begin{itemize}
  \item Hemiplegic gait.
  \item Paraplegic gait.
  \item Cerebellar ataxic gait.
  \item Sensory ataxic gait.
  \item Steppage gait.
  \item Dystrophic gait.
  \item Parkinsonian gait.
  \item Choreic gait.
  \item Apraxic gait.
  \item Antalgic gait.
  \end{itemize}
\end{itemize}
\end{latin}

\endgroup
\ifNEStandalone\expandafter\enddocument\fi
